\chapter{Introducció}
\label{chap:introduccio}

La recerca en intel·ligència artificial aplicada a jocs ha passat dels clàssics d'informació perfecta ---escacs, Go--- a entorns on l'estat no és tot visible i el resultat d'una acció no és determinista. Pokémon, simulat amb precisió de motor a Pokémon Showdown~\cite{pokemon_showdown} i exposat a Python via \textit{poke-env}~\cite{poke_env}, concentra aquests dos ingredients: informació oculta (moviments, objecte, habilitat i banc del rival) i estocasticitat de combat (rolls de dany, precisió, crítics, empats de velocitat). El problema no és un MDP clàssic observable; és un \textbf{joc estocàstic d'observació parcial} (POSG), proper a un POMDP~\cite{pomdp,kaelbling_pomdp} de dos jugadors a suma zero~\cite{von_neumann}.

Aquest Treball de Final de Màster és un estudi de \textbf{comparació de paradigmes}, no un intent de «resoldre Pokémon» ni de construir l'agent més fort possible contra humans. El format és \textbf{exclusivament Singles} \texttt{gen9randombattle}. El competitiu oficial de Dobles (VGC) queda fora d'abast ---es va eliminar del pla del projecte--- i no ocupa cap capítol d'aquesta memòria. Un sol format, un sol simulador, la mateixa matriu d'oponents: la diferència entre agents s'ha d'atribuir a la política, no a un canvi de metagame.

\section{Pregunta de recerca}

La pregunta operativa és:

\begin{quote}
\textit{En un POSG com \texttt{gen9randombattle}, quin paradigma de decisió s'acosta més al sostre de coneixement d'una heurística forta: regles explícites, cerca adversària d'un torn, cerca de Monte Carlo a diversos torns, o clonatge d'experts humans?}
\end{quote}

«Sostre de coneixement» no vol dir «el millor agent concebible». Vol dir el millor agent \emph{d'aquesta} escala heurística, mesurat a la mateixa bateria bot-contra-bot. El gauntlet de 28 agents respon amb un rànquing, no amb una demostració d'optimalitat.

Quatre paradigmes estan mesurats de cap a cap:

\begin{enumerate}
    \item \textbf{Heurístiques \texttt{v1}--\texttt{v14}}: ablació per construcció. Cada versió afegeix una peça que un expert humà anomenaria (dany, hazards, Tera, predicció de set, Yomi).
    \item \textbf{Minimax 1-ply \texttt{v15}--\texttt{v17}}: maximin analític sobre la resposta predita del rival, sense simulador intern.
    \item \textbf{Information-Set MCTS \texttt{v18}--\texttt{v20}}: 100 determinacions $\times$ 5 torns de \texttt{LocalSim}.
    \item \textbf{Imitació \texttt{v21}--\texttt{v22}}: el mateix classificador XGBoost de macro (mou vs canvia), dues piles d'execució (híbrid amb \texttt{v14} vs clon pur).
\end{enumerate}

El cinquè paradigma, \textbf{PPO / MaskablePPO} (\texttt{v23+}), té framework al repositori (\texttt{src/p05\_ppo\_drl/}) i \textbf{no} entra en aquesta matriu. S'avaluarà a part, gairebé segur \emph{només contra \texttt{v12}} ---el sostre bot-contra-bot---, no contra els 28 agents. Aquesta memòria el tracta com a treball futur.

\section{Objectius}

\begin{itemize}
    \item Formalitzar \texttt{gen9randombattle} com a POSG i justificar per què la mida del Pokédex \emph{no} entra a la variància d'una taxa de victòries binomial.
    \item Construir i documentar una infraestructura de simulació paral·lela (Showdown local + workers aïllats) capaç de 10.000 partides per cel·la.
    \item Desenvolupar l'escala heurística \texttt{v1}--\texttt{v14} com a ablació de coneixement, no com a catorze bots intercanviables.
    \item Implementar cerca 1-ply i IS-MCTS \emph{a sobre} de \texttt{v14}, amb telemetria de si la cerca s'executa i de si discrepa de l'expert.
    \item Entrenar un macro d'imitació sobre replays humans Elo~$1800+$ del \emph{mateix} format, i contrastar híbrid vs clon pur.
    \item Publicar una matriu $28\times 28$ amb protocol honest: $n=10.000$ per defecte, $n=1.000$ quan intervé MCTS, ladder humana de \texttt{v14} com a validació ecològica (no com a substitut del round-robin).
\end{itemize}

\section{Contribucions}

\begin{enumerate}
    \item \textbf{Protocol estadístic explícit.} Cada partida és un Bernoulli. Es deriva l'IC del 95\%, es justifica $n=10.000$ com a nivell de tesi ($\pm 0{,}98$~pp) i $n=1.000$ com a nivell de validació ($\pm 3{,}1$~pp), i es documenta per què MCTS no puja a 10k. El nombre de Pokémon de la generació no hi entra.
    \item \textbf{Escala d'ablació heurística amb tres salts i una inversió.} El dany extra no guanya Random Battles (\texttt{v1}--\texttt{v6} $\approx 44$--$46\%$). La posició (\texttt{v7}) i el tempo (\texttt{v9}) sí. Les mecàniques natives de Gen~9 (\texttt{v12}: Tera, lead de preview, switch-in) són el salt gros: 69,01\% global, 59,91\% vs Abyssal. La genealogia deia $\texttt{v14}\supset\texttt{v13}\supset\texttt{v12}$; el gauntlet diu el contrari. Això no és un bug: \texttt{v14} està dissenyat per humans.
    \item \textbf{Mesura de la cerca com a residual, no com a política.} Tant el minimax com l'MCTS s'executen en $\sim$16--18\% dels torns (la resta són dreceres de KO). La taxa de \emph{override} de \texttt{v14} predice qui guanya: 67\% / 71\% als no constrenyits, 37\% / 19\% als híbrids.
    \item \textbf{Contrast híbrid vs pur en imitació.} El mateix model, el mateix llindar $\tau=0{,}5525$. Conservar \texttt{v14} dona un jugador de nivell \texttt{v9}/\texttt{v11} (58,71\%). Treure'l cau per sota de \texttt{v1} (33,49\%).
    \item \textbf{Moral jeràrquica transversal.} Apareix tres vegades. El residual que substitueix l'expert perd; el que el conserva és l'única millora; cap no supera \texttt{v12}.
\end{enumerate}

No és una contribució d'aquesta passada presentar PPO com a experiment acabat, ni presentar Dobles/VGC com a part del corpus, ni citar 98 partides / 40,8\% / Elo~1151 com a mostra final de ladder.

\section{Estructura de la memòria}

\begin{itemize}
    \item El \textbf{Capítol~\ref{chap:marc_teoric}} situa el POSG, les mecàniques de Singles que defineixen $S$, $A$ i $\Omega$, i els quatre paradigmes (heurística, minimax, IS-MCTS, imitació). El RL/PPO s'hi explica com a rerefons del treball futur, no com a resultat.
    \item El \textbf{Capítol~\ref{chap:metodologia}} descriu Showdown, poke-env, la matriu $28\times 28$, i ---amb fórmula i taula--- per què 10.000 i 1.000 partides no són intercanviables.
    \item El \textbf{Capítol~\ref{chap:infraestructura}} documenta el maquinari i el flux de treball remot.
    \item El \textbf{Capítol~\ref{chap:implementacio}} detalla l'orquestrador master--worker, l'escala \texttt{v1}--\texttt{v14}, el maximin 1-ply, l'IS-MCTS i els dos agents IL.
    \item El \textbf{Capítol~\ref{chap:resultats}} és el cos empíric: protocol, escala heurística, inversió, minimax, MCTS, IL, heatmap, Elo, ladder, discussió.
    \item El \textbf{Capítol~\ref{chap:conclusions}} resumeix el que es pot afirmar i el que queda per PPO vs \texttt{v12}.
\end{itemize}
